% TOP_PROBLEM: {"problemId": "problem.smallest-number-of-variables-for-undecidable-hilbert-s-tenth-problem-over-z", "canonicalTitle": "Smallest number of variables for undecidable Hilbert's Tenth Problem over Z", "releaseRank": 416, "primaryDomain": "logic_foundations_set_theory", "primaryDomainLabel": "Logic, foundations & set theory", "displayStatus": "open_with_solved_subcases", "releaseStatus": "open_with_solved_subcases"}
% !TeX program = lualatex
% ATTEMPT_STATUS: unresolved
% Research continuation by GPT_6_Astra_Ultra, 27 September 2026.
\documentclass[11pt]{article}
\usepackage[margin=25mm]{geometry}
\usepackage{fontspec}
\setmainfont{Latin Modern Roman}
\usepackage{amsmath,amssymb,mathtools}
\usepackage{unicode-math}
\setmathfont{Latin Modern Math}
\usepackage{xurl,hyperref}
\hypersetup{colorlinks=true,urlcolor=blue}
\setcounter{secnumdepth}{0}
\setlength{\parindent}{0pt}
\setlength{\parskip}{5pt}
\setlength{\emergencystretch}{3em}
\begin{document}
\section*{\texorpdfstring{Smallest number of variables for undecidable Hilbert's Tenth Problem over Z}{Smallest number of variables for undecidable Hilbert's Tenth Problem over Z}}\label{416}
\subsection{Definitions and mathematical statement}
For a fixed nonnegative integer $n$, let $\mathrm{H10}_{{\mathbb{Z}}}(n)$ be the decision problem whose input is an arbitrary finite integer-coefficient polynomial in at most $n$ variables and whose yes-instances are those with an integer zero:
\[
 f\in{\mathbb{Z}}[x_1,\ldots,x_n],\qquad
 \exists (a_1,\ldots,a_n)\in{\mathbb{Z}}^n:\ f(a_1,\ldots,a_n)=0.
\]
Determine the least $n$ for which this problem is undecidable, meaning no algorithm halts with the correct answer on every input. Coefficients and degree are unrestricted and part of the finite input. The number of variables is the parameter being minimized; fixing degree instead, allowing only nonnegative solutions, or replacing ${\mathbb{Z}}$ by ${\mathbb{Q}}$ changes the problem. To identify the threshold, both lower-variable decidability and undecidability at that threshold are needed.
\par\smallskip\textit{Formulation and concept references:} \href{https://www.cs.umd.edu/users/gasarch/open/hilbert10openproblems.pdf}{[S1]}.

\subsection{Short English statement}
What is the smallest fixed number of variables for which deciding existence of an integer root of an arbitrary integer polynomial becomes algorithmically impossible?
\subsection{Research attempt}
\textit{GPT-6 Astra Ultra; 27 September 2026. Status: UNRESOLVED.}

This is a threshold question, so an undecidability construction in a
small number of variables is only one side of the answer. I prove the
elementary lower bound, formulate a precise height-bound obstruction,
and audit the variable accounting in several natural reductions. A
recent September 2026 manuscript must also be distinguished from the
older bound quoted in the formulation source.

\paragraph{The zero- and one-variable problems are decidable.}
With no variables the input is a constant integer and the answer is
whether it equals zero. For one variable, first handle the zero
polynomial and nonzero constants. Write a remaining polynomial as
$x^r g(x)$ with $g(0)\ne0$. If $r>0$, zero is already a root.
Otherwise, if $a\in\mathbb Z$ is a root, the equation
\[
 g(a)=c_m a^m+\cdots+c_1a+c_0=0
\]
shows that $a$ divides the nonzero integer $c_0$. There are finitely
many signed divisors, and evaluating at all of them decides whether
an integer root exists. The divisor enumeration can be slow but
terminates, which is the relevant criterion here. Consequently the
least undecidable variable count $n_{\min}$ satisfies $n_{\min}\ge2$.

For comparison, linear equations in any finite number of variables
are decidable. The equation $a_1x_1+\cdots+a_nx_n=b$ is soluble
exactly when $\gcd(a_1,\ldots,a_n)$ divides $b$, with the all-zero
coefficient case treated separately. Necessity follows by divisibility
and sufficiency by the extended Euclidean algorithm. This does not
settle $\mathrm{H10}_{\mathbb Z}(n)$ for that $n$, because its input
degree is unrestricted. Degree bounds and variable bounds cannot be
interchanged in a threshold claim.

\paragraph{Decidability is equivalent to a computable witness bound.}
Fix $n$. Integer solubility is always semidecidable: enumerate the
boxes $[-H,H]^n$ for $H=0,1,2,\ldots$ and evaluate the polynomial
exactly. A solution eventually appears if one exists. The algorithm
may run forever otherwise.

There is a total decision algorithm for this fixed $n$ if and only if
there is a total computable function $B(f)$ such that every soluble
input has a zero $a$ with
\[
 \max_i|a_i|\le B(f).                                 \tag{1}
\]
A bound gives an algorithm by searching the finite box. Conversely,
given a decision algorithm, return zero on a negative input. On a
positive input, enumerate until a witness is found and return its
height. The promised positive answer ensures termination of this
second phase. This proves the equivalence without any complexity
claim or assumption that a minimal solution is easy to find.

Therefore an undecidability proof rules out all computable bounds of
the form (1), not just polynomial or exponential bounds. Very large
solutions in a finite sample do not establish undecidability: they
can at most refute a particular proposed numerical bound. Conversely,
searching large boxes without finding a root gives no negative
certificate unless a valid bound has already been proved.

\paragraph{Variable accounting for elementary encodings.}
Over the integers a finite system $f_1(x)=\cdots=f_r(x)=0$ is
equivalent to the single equation
\[
 f_1(x)^2+\cdots+f_r(x)^2=0.                           \tag{2}
\]
The proof uses nonnegativity of integer squares. This step adds no
variables, but can double the degree. It does not eliminate any
auxiliary variables used to define the $f_i$ in the first place.
The same encoding over $\mathbb C$ fails, for example at $(1,i)$
in $U^2+V^2=0$; a field change is not innocuous.

To convert $r$ nonnegative unknowns into integer unknowns, Lagrange's
four-square theorem permits
\[
 x_i=y_{i1}^2+y_{i2}^2+y_{i3}^2+y_{i4}^2.
                                                        \tag{3}
\]
Every right side is nonnegative, and every nonnegative integer has
such a representation. Thus the direct substitution replaces $r$
unknowns by $4r$ and at most doubles the degree. This is a valid but
not optimal conversion. In particular, a nine-variable theorem over
the nonnegative integers does not become a nine-variable theorem over
$\mathbb Z$ by changing the domain label. More efficient encodings
must be proved and their witnesses counted explicitly.

There is a different bookkeeping issue in a universal polynomial
$U(a,z_1,\ldots,z_m)$. If the reduction receives an input integer
$a$ and substitutes it as a coefficient, the resulting equation has
$m$ existential unknowns, not $m+1$. If $a$ remains existentially
quantified, it is a different problem and the reduction loses the
original input. Both errors can shift a purported threshold.

\paragraph{Why a computable pairing does not give a free compression.}
Cantor pairing gives a computable bijection $\mathbb N^2\to\mathbb N$,
but its inverse coordinates are not polynomials. One elementary
obstruction is as follows. Suppose polynomials $P,Q\in\mathbb Q[t]$
take nonnegative values on all sufficiently large integers and satisfy
\[
 \frac{(P(t)+Q(t))(P(t)+Q(t)+1)}2+Q(t)=t.              \tag{4}
\]
If both are constant, the left side is constant. Otherwise their
eventual nonnegativity ensures that the leading coefficients of the
nonconstant polynomials are positive, so $P+Q$ has degree
$d=\max(\deg P,\deg Q)\ge1$. The left side has degree $2d\ge2$,
contradicting (4). Thus this particular computable coding does not
provide polynomial decoding without further existential witnesses.
The argument does not exclude all ingenious Diophantine encodings;
it pinpoints a failure in one common shortcut.

\paragraph{The current upper-bound source must be stated carefully.}
The older eleven-integer-variable theorem is recorded in [S1] and
proved in Sun's earlier work [E1]. A newer manuscript, Sun's
\emph{Ten unknowns for Hilbert's tenth problem over the integers},
arXiv:2609.01594v2, dated 2 September 2026, states a ten-variable
improvement. Its Theorem 1.1 represents every recursively enumerable
set by $Q_{\mathcal A}(a,z_1,\ldots,z_{10})$ with $a$ a specialized
parameter; Corollary 1.1 states undecidability with ten integer
unknowns. I checked that statement and its variable/domain convention,
but have not independently verified its full proof. Thus the
established older bracket is $2\le n_{\min}\le11$, and the recent
manuscript reports the improvement $n_{\min}\le10$. Even accepting
that improvement does not identify the minimum.

\paragraph{Verification and outcome.}
Exact finite tests check the signed-divisor root algorithm against
bounded exhaustive evaluation on small polynomials, the integer
sum-of-squares equivalence, and four-square coverage on a stated finite
range. Those tests verify implementation examples, not the cited
undecidability results. The first unresolved step toward the exact
threshold is either a decision method for further fixed variable
counts or a smaller-variable undecidability reduction with complete
witness accounting. The proved elementary lower bound and equivalence
(1) do not bridge that gap. The requested minimum remains
\textbf{UNRESOLVED}.

\subsection{Sources}
\begin{itemize}
\item[S1]William Gasarch, Hilbert's Tenth Problem: Refinements and Variants (2021), Notation 3.1 and Theorems 4.1-4.2.
\url{https://www.cs.umd.edu/users/gasarch/open/hilbert10openproblems.pdf}.
\textit{Formulation source.}
\item[E1]Zhi-Wei Sun, \emph{Further results on Hilbert's Tenth Problem},
arXiv:1704.03504.
\url{https://arxiv.org/abs/1704.03504}.
\textit{Earlier eleven-integer-variable theorem; consulted alongside
the older bound in [S1] on 27 September 2026.}
\item[E2]Zhi-Wei Sun, \emph{Ten unknowns for Hilbert's tenth problem
over the integers}, arXiv:2609.01594v2 (2 September 2026),
Theorem 1.1 and Corollary 1.1.
\url{https://arxiv.org/html/2609.01594v2}.
\textit{Recent manuscript reporting a ten-variable improvement. The
statement, parameter count, and integer domain were checked; its entire
proof was not independently audited. Consulted 27 September 2026.}

\end{itemize}
\end{document}
